\section{The finite family of initial marks}
\label{sec:area_law_initial_mark_family}

The sparse shifts may be chosen in every distance layer. Their cell counts
have a summable geometric bound, so only finitely many layers contribute
belt cells. We can therefore form the finite union of their actual marked
points and assign to each point its smallest incident cell side; see
\cite[Section~11]{OpenAI2026AreaLaw}.

% Source: 10-geometry.tex, geometry:belt-count, lines 220--233;
% initial marks, lines 325--330; geometry:total-repairs, lines 668--692.
% Only the zero-polynomial summation for initial cells and marks is proved
% here. Point separation is the first assertion of geometry:initial-stars,
% lines 325--339, with its proof at lines 352--359. Repairs remain separate.

\begin{theorem}[Sparse shifts with finite support]
    \label{thm:area_law_finite_sparse_belt_family}
    \lean{TNLean.PEPS.AreaLaw.Geometry.exists_finitely_supported_sparse_belt_shifts}
    \leanok
    \uses{def:area_law_belt_cells,def:area_law_fine_layer_indices,
        def:area_law_fine_belt_scales}
    For every finite domain $\Lambda$, cut $A$, origin $o$ and
    nonnegative integer $C_0$, there are coordinate residues $a_k,b_k$
    at every scale. Write $F_k$ for their actual belt-cell index set.
    These choices satisfy
    \begin{align}
        |F_k| & \le64(2C_0+1)^2|\partial_\Lambda A|\,2^{-\delta_0 k/2}
        \qquad(k\ge0),
    \end{align}
    and the set $\{k:F_k\ne\varnothing\}$ is finite.
\end{theorem}
\begin{proof}\leanok
    \uses{thm:area_law_sparse_dyadic_belts,thm:area_law_polynomial_summability}
    Choose a sparse pair of residues at each scale. The integer counts
    $|F_k|$ are dominated by a fixed multiple of a summable geometric
    sequence. Their real-valued series is therefore summable, so its
    terms tend to zero. Eventually $|F_k|<1$, which forces $F_k$ to be
    empty. Only finitely many indices remain.
\end{proof}

\begin{theorem}[Uniform count of all initial marks]
    \label{thm:area_law_uniform_initial_mark_count}
    \lean{TNLean.PEPS.AreaLaw.Geometry.exists_uniform_initial_mark_bound}
    \leanok
    \uses{def:area_law_belt_marks,def:area_law_belt_cells,
        def:area_law_fine_layer_indices,def:area_law_fine_belt_scales}
    There is a positive constant $B$ such that, for every finite domain,
    cut, origin, nonnegative integer $C_0$ and lower index $k_0$, one can
    choose the residues from the preceding theorem and a finite set $M$
    of actual marks with
    \begin{align}
        x\in M
            & \quad\Longleftrightarrow\quad
        \exists k\ge k_0,\ x\in M_{\ell_k}(F_k),   \\
        |M| & \le B(2C_0+1)^2|\partial_\Lambda A|.
    \end{align}
    The chosen residues retain their per-layer cell estimate and finite
    support. The constant $B$ precedes all the geometric data; the
    support cutoff may depend on those data.
\end{theorem}
\begin{proof}\leanok
    \uses{thm:area_law_finite_sparse_belt_family,thm:area_law_belt_mark_count,
        thm:area_law_polynomial_finite_sums}
    Let $B_0>0$ be the uniform finite-sum budget for the zero polynomial
    exponent, and take $B=576B_0$. Form the union of the actual marks
    over the finite support restricted to $k\ge k_0$. Its membership is
    exactly the stated condition. A finite union has at most the sum of
    the individual cardinalities, and each cell contributes at most
    nine marks. Thus its count is at most
    \begin{align}
        9\sum_{k\ge k_0}|F_k|
         & \le576(2C_0+1)^2|\partial_\Lambda A|
        \sum_{k\in H}2^{-\delta_0 k/2}
        \le B(2C_0+1)^2|\partial_\Lambda A|,
    \end{align}
    where $H$ is the finite retained support.
\end{proof}


\begin{theorem}[Separation of fine-layer marks]
    \label{thm:area_law_fine_layer_mark_separation}
    \lean{TNLean.PEPS.AreaLaw.Geometry.fineLayer_marks_dist_ge}
    \leanok
    \uses{def:area_law_belt_marks,def:area_law_fine_layer_indices,
        def:area_law_fine_belt_scales}
    Fix an origin $o$ and a finite endpoint set $Z$, and form all layers
    from these data with $C_0\ge2$. Let $k,h\ge50{,}000{,}000$. If $x$ and $y$ are
    distinct marks of actual fine-layer cells in layers $k$ and $h$,
    respectively, then
    \begin{align}
        \operatorname{dist}_\infty(x,y)\ge\frac14\max(t_k,t_h),
        \qquad t_j=2^{\ell_j}.
    \end{align}
    Cell-boundary marks are included, and the estimate applies to every
    fine-layer cell, hence to every selected belt cell.
\end{theorem}
\begin{proof}\leanok
    \uses{thm:area_law_local_layer_indices,thm:area_law_adjacent_fine_meshes,
        thm:area_law_belt_marks_quarter_mesh,thm:area_law_affine_mesh_point_separation,
        thm:area_law_fine_cell_containment}
    Suppose that the distance is less than $t_k/4$. Each mark is in the
    closure of its actual layer. Layer locality forces $k\le h+1$.
    Monotonicity and the adjacent-scale estimate give
    $\ell_k\le\ell_h+1$, so both marks lie on the translated mesh of
    spacing $t_k/4$. Distinct mesh points cannot be closer than this
    spacing. Reversing $k$ and $h$ proves both lower bounds and hence
    their maximum.
\end{proof}

\begin{theorem}[Counted and separated initial marks]
    \label{thm:area_law_uniform_separated_initial_marks}
    \lean{TNLean.PEPS.AreaLaw.Geometry.exists_uniform_separated_initial_marks}
    \leanok
    \uses{thm:area_law_uniform_initial_mark_count,def:area_law_belt_marks}
    There is $B>0$ such that, for every finite domain, cut, origin,
    integer $C_0\ge2$ and lower index $k_0\ge50{,}000{,}000$, there are
    sparse residues with the per-layer estimates and finite support,
    the exact finite union $M$ of their actual marks above $k_0$, and a
    positive side $s(v)$ for each $v\in M$, satisfying
    \begin{align}
        |M| & \le B(2C_0+1)^2|\partial_\Lambda A|,     \\
        \operatorname{dist}_\infty(v,w)
            & \ge\frac14\max(s(v),s(w))\qquad(v\ne w).
    \end{align}
    For each $v$, some selected belt cell in a layer $k\ge k_0$ has $v$
    as an actual mark and side $s(v)=t_k$. Every selected belt cell in a layer $h\ge k_0$ that has $v$ as one
    of its actual marks has side at least $s(v)$.
\end{theorem}
\begin{proof}\leanok
    \uses{thm:area_law_uniform_initial_mark_count,
        thm:area_law_fine_layer_mark_separation,thm:area_law_adjacent_fine_meshes}
    Choose the counted finite union and residues from the preceding
    construction. For each point of $M$, take the least layer index
    admitting an incident selected cell. The fine-cell exponent is
    nondecreasing, so this gives a positive smallest incident side.
    Apply fine-layer mark separation to the two witnessing cells.
    When $M$ is empty the assignment and pairwise condition are empty.
\end{proof}
